% Shared threat model for the three privacy experiments.
\subsection{Adversary and Observables}
\label{sec:eval-threat}

We consider the passive public-chain observer defined in Section~\ref{sec:problem}.
This adversary has unrestricted read access to the ledger and observes all transactions and events generated by TPS-ZKPI, but has no direct access to the plaintext policy or the witness.
Under the assumptions stated in Section~\ref{subsec:trust-model}, the zero-knowledge proof reveals no information beyond the public statement, and the hiding commitments conceal the committed policy.
These cryptographic protections do not, by themselves, prevent policy information from leaking through observable metadata.

We consider three adversary goals, one per privacy objective:
(i)~recover the policy's \emph{content} (PCH, Section~\ref{sec:eval-pch});
(ii)~distinguish two \emph{semantically related} policies (PI, Section~\ref{sec:eval-pi}); and
(iii)~recover the policy's \emph{structure or parameter sizes} (SPP, Section~\ref{sec:eval-spp}).

To evaluate attacks against these goals, we extract the following metadata features from each policy execution, consisting of a commit transaction followed by an access transaction:
%
\begin{itemize}
	\item \textbf{Gas:} \texttt{commitGasUsed}, \texttt{accessGasUsed}.
	\item \textbf{Calldata length:} \texttt{commitCalldataLength}, \texttt{accessCalldataLength}.
	\item \textbf{Proof size:} \texttt{proofSize}.
	\item \textbf{Event logs:} \texttt{commitLogsCount}, \texttt{commitLogsSize},
	      \texttt{accessLogsCount}, \texttt{accessLogsSize}.
\end{itemize}
%
The fixed circuit and commitment-only interface keep calldata length, proof size, and event-log counts and sizes constant across the evaluated policies.
These features remain in the attack pipeline but provide no variation for the attacks to exploit.
Any policy-dependent variation in the selected feature set must therefore occur in the two gas measurements (\texttt{commitGasUsed} and \texttt{accessGasUsed}).
